\documentclass[12pt]{article}
\usepackage[T1]{fontenc}
\usepackage[utf8]{inputenc}
\usepackage{lmodern}
\usepackage{textcomp}
\usepackage{enumitem}
\usepackage{geometry}
\usepackage{graphicx}
\usepackage{amsmath, amssymb}
\geometry{margin=1in}
\begin{document}
\begin{center}
History - K-12 Level Assessment 6 \
Total Marks: 40 \
Answer all questions.
\end{center}

\section*{Multiple Choice (10 marks)}
\begin{enumerate}[label=\arabic*.]
\item The Egyptian system of hieroglyphics:
\begin{enumerate}[label=(\alph*)]
    \item did not use pictographs
    \item appears to have developed suddenly
    \item was the earliest form of writing in the world
    \item all of the above
\end{enumerate}
\item In comparison to the Middle Paleolithic, Upper Paleolithic tool technologies:
\begin{enumerate}[label=(\alph*)]
    \item emphasized flakes over blades.
    \item changed rapidly.
    \item involved elaborate preparation of core tools.
    \item all of the above
\end{enumerate}
\item What do anthropologists NOT need to know to determine when people entered the Americas from Siberia?
\begin{enumerate}[label=(\alph*)]
    \item When sled dogs were first domesticated
    \item When the Bering Strait was exposed and open for travel
    \item When eastern Siberia was first inhabited
    \item The age of the earliest New World sites
\end{enumerate}
\item In comparison to an Oldowan tool, a hand axe \_\_\_\_\_\_\_\_\_\_\_\_.
\begin{enumerate}[label=(\alph*)]
    \item is made out of stone, wood, or bone.
    \item needs constant cleaning and maintenance.
    \item requires metal and a source of heat to be forged
    \item requires many more flakes to be removed.
\end{enumerate}
\item Unlike most other early civilizations, Minoan culture shows little evidence of:
\begin{enumerate}[label=(\alph*)]
    \item trade.
    \item warfare.
    \item the development of a common religion.
    \item conspicuous consumption by elites.
\end{enumerate}
\end{enumerate}

\section*{True / False (10 marks)}
\textit{Answer True or False}
\begin{enumerate}[label=\arabic*.]
\setcounter{enumi}{5}
\item True or False: Having a forward-thrusting lower face is referred to as being prognathous.
\item True or False: At the end of the Miocene, many species of ape became extinct due to significant changes in the environment.
\item True or False: The peak of Minoan civilization was bracketed by an earthquake and a volcanic eruption.
\item True or False: Ancestral Pueblo elites gained power by storing food surpluses during good years for distribution in times of scarcity.
\item True or False: Mastabas are mud-brick structures built over elite tombs that later evolved into pyramids in ancient Egypt.
\end{enumerate}

\section*{Long Form Response (20 marks)}
\textit{Answer each question with a clear and complete explanation.}
\begin{enumerate}[label=\arabic*.]
\setcounter{enumi}{10}
\item Discuss the agricultural practices of the Hassunan and Samarran sites in Mesopotamia, focusing on their similarities and the impact on their societies.
\item Discuss the underlying religious beliefs of the Olmec civilization and explain how these beliefs, as reflected in their art and iconography, connected them across different territories.
\end{enumerate}

\vfill
\noindent\textit{End of Paper}
\end{document}